% Generated draft: requires mathematical and source review.
\documentclass[11pt]{article}
\usepackage[margin=25mm]{geometry}
\usepackage{fontspec}
\setmainfont{Latin Modern Roman}
\usepackage{amsmath,amssymb,mathtools,unicode-math}
\setmathfont{Latin Modern Math}
\usepackage{enumitem,needspace,xurl,hyperref}
\setlength{\parindent}{0pt}
\setlength{\parskip}{5pt}
\setlength{\emergencystretch}{3em}
\newcommand{\N}{\mathbb{N}}
\newcommand{\Z}{\mathbb{Z}}
\newcommand{\Q}{\mathbb{Q}}
\newcommand{\R}{\mathbb{R}}
\newcommand{\C}{\mathbb{C}}
\newcommand{\F}{\mathbb{F}}
\newcommand{\A}{\mathbb{A}}
\newcommand{\PP}{\mathbb{P}}
\newcommand{\E}{\mathbb{E}}
\newcommand{\eps}{\varepsilon}
\newcommand{\dd}{\,\mathrm{d}}
\begin{document}
\section*{UnsolvedMath ID: 801022}
\textbf{Source problem code:} catalogue entry
\subsection{Definitions and mathematical statement}
A group \(G\) is \emph{finitely generated} if some finite subset generates \(G\). It is \emph{residually finite} if
\[
\forall g\in G\setminus\{1_G\}\;\exists\text{ a finite group }F\;\exists\text{ a homomorphism }\varphi:G\to F\quad \varphi(g)\ne 1_F.
\]
For a group \(G\) with a specified finite generating tuple \(S=(s_1,\ldots,s_d)\), its word problem is the decision problem for
\[
W(G,S)=\{w\in\{s_1^{\pm1},\ldots,s_d^{\pm1}\}^{*}:w=_G1_G\}.
\]
An unsolvable word problem means that this set is undecidable. For groups outside this finitely generated setting, the effective encoding of generators and words must also be specified.

The supplied record contains three questions. It does not define a \(\Pi\)-group, effective residual finiteness, or \(S_{\mathrm{rec}}\), and it omits the antecedent of ``such a group \(L\).'' Consequently, the following is a conditional transcription of the intended questions, rather than a complete standalone formulation. Write \(\mathcal P\) for the source's unspecified class of \(\Pi\)-groups, \(\mathcal E\) for its unspecified class of effectively residually finite groups, and \(\mathcal H\) for the missing hypotheses on \(L\). These symbols retain missing source conditions; they do not supply definitions for them.
\begin{enumerate}
\item Does there exist a group \(L\) satisfying \(\mathcal H\) whose word problem is unsolvable, using the generating data and effective encoding intended by the source? The hypotheses \(\mathcal H\) and these effective data are absent from the supplied record.
\item Does there exist a group \(G\) such that
\[
G\in\mathcal P,\qquad G\text{ is finitely generated},\qquad G\text{ is residually finite},\qquad G\notin\mathcal E?
\]
\item Does there exist a residually finite group \(G\in\mathcal P\) for which there is no injective homomorphism
\[
G\longrightarrow S_{\mathrm{rec}}?
\]
No finite-generation assumption is stated for this third question. The source attributes to Morozov a construction of a \(\Pi\)-group not isomorphic to any subgroup of \(S_{\mathrm{rec}}\), and asks whether a group with that nonembedding property can also be residually finite. It additionally states that such a group cannot be effectively residually finite. Both the attributed construction and this additional assertion are retained as supplied background, without independent verification; their precise content depends on the missing definitions.
\end{enumerate}
\subsection{Short English statement}
Problem 34 asks whether a previously described group \(L\) can have an undecidable word problem, whether a finitely generated residually finite \(\Pi\)-group can fail effective residual finiteness, and whether a residually finite \(\Pi\)-group can fail to embed into \(S_{\mathrm{rec}}\). The antecedent for \(L\) and several essential definitions are missing, so a complete standalone statement requires source clarification. The registry labels the problem open; that label is supplied metadata.
\subsection{Sources}
\begin{itemize}
\item Supplied local UnsolvedMath registry record, permanent ID \(801022\), exact source problem code \texttt{catalogue entry}; collection: AMR Open Problem Lists; title: \emph{Open Problems in Mathematical Logic --- Problem 34}. The record attributes the problem to Barmpalias, Gao et al., with proposed year 2026, and labels it open. These are supplied attribution and status fields, not independently verified conclusions.
\item Barmpalias, Gao et al., \emph{Open Problems in Mathematical Logic} (2026), Problem 34. Supplied source URL: \url{https://arxiv.org/abs/2608.26628}. The background records access on 2026-10-16 and extraction by \texttt{arxiv\_latex}; this is supplied provenance, not a claim of independent access or verification.
\item Morozov, \emph{Higman's question revisited} (2000), as cited in the supplied statement for the construction of a \(\Pi\)-group not isomorphic to a subgroup of \(S_{\mathrm{rec}}\). No fuller bibliographic details or URL are supplied for this reference, and the attribution has not been independently verified.
\end{itemize}
\end{document}
